\section{Pre-Christian Continuities Received and Corrected by the Fathers}
\label{sec:pre-christian-continuities}

When St.~Augustine confronted the Platonist philosophers in the ninth book of \emph{The City of God}, he refused to quarrel over a bare word once the substance of the faith stood clear. Referring to the created immortal intelligences whom Plato's \emph{Timaeus} calls ``gods,'' Augustine wrote:
\begin{quotation}
If the Platonists prefer to call these angels gods rather than demons, and to reckon them with those whom Plato, their founder and master, maintains were created by the supreme God, they are welcome to do so, for I will not spend strength in fighting about words. For if they say that these beings are immortal, and yet created by the supreme God, blessed but by cleaving to their Creator and not by their own power, they say what we say, whatever name they call these beings by. (St.~Augustine, \emph{City of God} IX.23)
\end{quotation}
That passage captures the patristic posture toward the pre-Christian world. The Fathers neither cut Catholic angelology off from the intellectual inheritance of antiquity nor surrendered revealed doctrine to Hellenistic cosmology. Across four distinct streams---the Greek Septuagint, extra-canonical Second Temple Jewish writings, Philo of Alexandria, and the Greek poets and philosophers from Hesiod to the Neoplatonists---the Church sifted what was true, named what was apocryphal, and corrected whatever obscured the sole creatorship of God and the sole mediation of Christ.

\subsection{The Septuagint Bridge}
\label{subsec:lxx-bridge}

The earliest and deepest transformation happened inside Judaism itself when the Alexandrian translators rendered the Hebrew Scriptures into Greek. Hebrew vocabulary for heavenly beings was varied and fluid: \gk{mal'akim} (``messengers''), \gk{bene ha-elohim} (``sons of God''), \gk{elohim} (``gods'' or ``mighty ones''), and \gk{elilim} (``idols'' or ``nothings''). The Septuagint forged from these words the precise Greek terminology---\gk{angeloi}, \gk{archontes}, and \gk{daimonia}---that passed directly into the New Testament and the teaching of the Fathers. Eight Septuagint passages left a lasting mark on Catholic angelology.

In \textbf{Deuteronomy~32:8--9}, where the Hebrew Masoretic text and the Clementine Vulgate read that the Most High set the bounds of the nations ``according to the number of the children of Israel'' (\latin{juxta numerum filiorum Israel}), the Septuagint reads: ``When the Most High divided the nations, when he separated the sons of Adam, he set the bounds of the nations according to the number of the angels of God'' (\gk{angel\={o}n theou}). Origen, St.~Basil, St.~John Chrysostom, and St.~Jerome (\emph{Commentary on Daniel}~10) joined this verse to Daniel~10:13 and 20--21 to ground the common patristic teaching on national guardian angels. Pseudo-Dionysius (\CH{9.2--4}) cited the Septuagint wording explicitly while guarding it against any polytheistic or Gnostic distortion: God did not partition the earth with independent angels or take Israel by lot, but one universal Providence placed holy angels over every nation to draw willing peoples toward the one true God, and Gentile idolatry sprang from human self-will rather than angelic unfaithfulness.

In \textbf{Psalm~8:6} [Heb~8:6; Eng.~8:5], the Septuagint translated Hebrew \gk{elohim} as \gk{angelous}: ``Thou madest him a little less than angels, thou hast crowned him with glory and honour.'' The Latin Vulgate (\latin{paulominus ab angelis}) and the Epistle to the Hebrews (Heb~2:7) adopted this phrasing directly. Through it, patristic theology articulated both the natural order of rational creation---God above angels, angels above man---and the condescension of the Incarnation, in which the Son accepted a human nature temporarily lower than the angels before exalting humanity above every heavenly rank.

In \textbf{Psalm~95:5} [96:5], the Septuagint enacted a decisive break with Greek philosophical usage by rendering the Hebrew word for worthless idols as \gk{daimonia}: ``For all the gods of the heathen are devils [\gk{daimonia}]: but the Lord made the heavens.'' In classical Greek, \gk{daim\={o}n} could signify a neutral or benevolent intermediary spirit. By fixing \gk{daimonia} upon the usurping spirits behind pagan idolatry, the Septuagint drew the line that St.~Justin, Athenagoras, and St.~Augustine (\emph{City of God} IX.19) maintained uniformly: in Christian Scripture, \gk{daim\={o}n} never designates a holy angel, but always a wicked spirit.

In \textbf{Psalm~96:7} [97:7], the Septuagint translated ``all gods'' as ``all his angels'': ``Let all that worship graven images be ashamed, who boast of their idols; worship him, all ye his angels.'' Adopted in Hebrews~1:6 (``And let all the angels of God adore him''), this verse became a primary anti-Arian proof in St.~Athanasius (\emph{Against the Arians}), St.~Ambrose (\emph{On the Faith}), and St.~Gregory of Nazianzus: the Son is uncreated God adored by every angel, never a created spirit numbered among them.

In \textbf{Psalm~103:4} [104:4], the Septuagint rendered ``Who makes his angels spirits [\gk{pneumata}], and his ministers a flaming fire [\gk{pyr phlegon}],'' quoted verbatim in Hebrews~1:7. St.~Basil (\emph{On the Holy Spirit} XVI.38), St.~Gregory of Nazianzus (\emph{Oration} 28.31), St.~Augustine (\emph{City of God} XV.23), and St.~Gregory the Great (\emph{Homily~34 on the Gospels}~\S 8) turned repeatedly to this verse to distinguish what an angel is by nature (\latin{spiritus}, ``spirit'') from what an angel is by office (\latin{angelus}, ``messenger''), and to weigh how Scripture speaks of angelic subtleness and fiery ardor.

In \textbf{Isaias~9:5 LXX} [Vulgate/Heb~9:6], the Greek Bible gave the promised Child a title absent from the Latin Vulgate: ``and his name is called the Messenger of great counsel'' (\gk{Megal\={e}s boul\={e}s angelos}). In the second century, St.~Justin Martyr (\emph{Dialogue with Trypho} 56--61, 128) and St.~Irenaeus invoked this title to identify the pre-incarnate Word in Old Testament appearances of the Angel of the Lord. When fourth-century Arians seized on the word \gk{angelos} to demote the Son to a created angel, post-Nicene Fathers---Athanasius, Augustine (\emph{On the Trinity}~III), and Pseudo-Dionysius---clarified the title strictly by mission rather than nature. As Dionysius explains (\CH{4.4}), ``even Jesus Himself, when He had come to manifest the good work of our beneficent salvation, was called Angel of Great Counsel. For, as He Himself says, after the manner of an Angel, `Whatsoever He heard from the Father, He announced to us.'''

In \textbf{Daniel~10:13 and 20--21}, the Greek version of Theodotion received by the Church names ``the prince of the kingdom of the Persians,'' ``the prince of the Greeks,'' and ``Michael, one of the princes... Michael your prince.'' St.~Jerome (\emph{Commentary on Daniel}~10), St.~Gregory the Great (\emph{Moralia in Job}~XVII), and Pseudo-Dionysius (\CH{9.2--4}) debated whether Persia's resisting prince was a good national angel pleading for his province or an evil spirit, while agreeing that Michael stands over Israel and the Church under God's single providence.

Finally, in \textbf{Job~1:6 and 38:7}, the Septuagint translated Hebrew ``sons of God'' directly as ``angels of God'': ``behold, the angels of God came to stand before the Lord, and the devil came with them'' (Job~1:6), and ``When the stars were made, all my angels praised me with a loud voice'' (Job~38:7). (By contrast, at Genesis~6:2 the standard Vaticanus Septuagint retains ``sons of God,'' though Codex Alexandrinus and several early writers read ``angels of God.'') Job~38:7 LXX became a cornerstone in St.~Basil (\emph{Hexaemeron}~I.5), St.~Ambrose (\emph{Hexameron}~I.5), St.~Jerome, and St.~Augustine (\emph{City of God} XI.9), proving that the angels already existed when the stars were fashioned and framing the patristic debate between pre-cosmic creation in the East and simultaneous Day-One creation in the West.

\subsection{Extra-Canonical Jewish Writings (Apocryphal)}
\label{subsec:apocrypha}

Between the third century before Christ and the first Christian century, Jewish apocalypses and testaments expanded biblical hints into elaborate narratives of angelic watches, heavenly tiers, and named archangels. The Church never received these writings into the canon of Sacred Scripture; they stand graded as \textbf{Apocryphal}. Yet early Christian authors knew them well, and the patristic sifting of three texts---\emph{1~Enoch}, the \emph{Book of Jubilees}, and the \emph{Testaments of the Twelve Patriarchs}---settled lasting boundaries of Catholic angelology.

\textbf{\emph{1~Enoch}} (Apocryphal) elaborates Genesis~6:1--4 into the myth of the Watchers---angels who lusted after mortal women, descended to earth, taught forbidden arts, and begot giants whose disembodied spirits became demons---and in chapter~20 lists seven watching archangels by name. Its renown in early Christianity grew from the canonical Epistle of St.~Jude, which quotes an Enochic judgment oracle:
\begin{quotation}
Now of these Enoch also, the seventh from Adam, prophesied, saying: Behold, the Lord cometh with thousands of his saints: To execute judgment upon all and to reprove all the ungodly for all the works of their ungodliness, whereby they have done ungodly: and for all the hard things which ungodly sinners have spoken against God. (Jude~1:14--15)
\end{quotation}
In the second century, St.~Justin Martyr received the Watchers explanation of Genesis~6 and turned it directly against Greek mythology in his \emph{Second Apology}~(5):
\begin{quotation}
God... committed the care of men and of all things under heaven to angels whom He appointed over them. But the angels transgressed this appointment, and were captivated by love of women, and begat children who are those that are called demons. (St.~Justin Martyr, \emph{Second Apology}~5)
\end{quotation}
Justin argued that pagan poets mistook those fallen angels and their demonic brood for Olympian gods. By the fifth century, however, the mature patristic tradition rejected both the Watchers fable and any canonical standing for the Enochic compilation. In \emph{The City of God}~(XV.23), St.~Augustine drew the decisive line:
\begin{quotation}
We cannot deny that Enoch, the seventh from Adam, left some divine writings, for this is asserted by the Apostle Jude in his canonical epistle. But it is not without reason that these writings have no place in that canon of Scripture which was preserved in the temple of the Hebrew people by the diligence of successive priests... So that the writings which are produced under his name, and which contain these fables about the giants, saying that their fathers were not men, are properly judged by prudent men to be not genuine... all of which, after careful examination, have been set apart from canonical authority under the title of Apocrypha. (St.~Augustine, \emph{City of God} XV.23)
\end{quotation}
Augustine read Genesis~6:2 as the sons of Seth marrying the daughters of Cain, insisted that 2~Peter~2:4 refers to the primordial fall of the devil and his angels rather than carnal union, and declared that he could by no means believe God's holy angels had so fallen. The Church's liturgical discipline followed the same restraint: despite \emph{1~Enoch}'s catalogue of seven archangels, the Synod of Rome under Pope St.~Zacharias (745) and subsequent Roman discipline recognized only the three names attested in canonical Scripture---Michael, Gabriel, and Raphael.

The \textbf{\emph{Book of Jubilees}} (Apocryphal) reflects two themes familiar to early Christian readers: an account of angels created on the first day alongside light and heaven, and in chapter~10 a narrative in which destructive spirits afflict Noah's descendants, Noah prays for relief, and the chief spirit Mastema begs God to leave some spirits under his command, whereupon God permits one-tenth to remain while consigning nine-tenths to condemnation. Christian demonology shared the biblical conviction that evil spirits act solely within divine permission (already plain in Job~1--2 and Luke~8:31), yet rejected both the Watchers origin of demons and a quasi-official prosecutor like Mastema in favor of the single voluntary rebellion of Lucifer through pride. Likewise, when St.~Augustine (\emph{City of God} XI.9; \emph{Literal Commentary on Genesis}~I--IV) argued that angels were created on the first day under the name of ``light,'' he grounded that exegesis strictly in Genesis~1:1--5 rather than in \emph{Jubilees}.

The \textbf{\emph{Testaments of the Twelve Patriarchs}} (Apocryphal), which include the Testament of Reuben's seven spirits of error and apocalyptic tiers of heaven, drew occasional notice from Origen and Tertullian as an extra-canonical reflection of Jewish expectation and moral combat. Yet when St.~Cyril of Jerusalem, St.~Gregory the Great, and Pseudo-Dionysius (\CH{6}) ordered the nine angelic choirs, they bypassed multi-heaven apocryphal itineraries altogether and built strictly upon canonical Old Testament and Pauline titles (Is~6, Ezech~1 and~10, Eph~1:21, Col~1:16, 1~Thess~4:15 [4:16], Jude~1:9).

\subsection{Philo of Alexandria}
\label{subsec:philo}

In first-century Alexandria, the Jewish philosopher Philo sought to harmonize the Torah with Middle Platonism across treatises such as \emph{De gigantibus}, \emph{De somniis}, \emph{De confusione linguarum}, and \emph{De opificio mundi}. Reading Genesis~6:1--4, Jacob's ladder, the Tower of Babel, and the creation narrative allegorically, Philo identified the \gk{angeloi} of Moses with the incorporeal rational souls and divine \gk{logoi} that Greek philosophers called \gk{daimones}---pure intelligences traversing the air as ambassadors of God's thoughts and ministers of the divine Logos.

Alexandrian Christian Fathers, above all Clement of Alexandria and Origen, inherited Philo's spiritual exegesis and his clarity about angels as incorporeal rational ministers serving the Word. Yet Catholic theology corrected three Philonic tendencies that proved incompatible with the Gospel. First, the Fathers refused Philo's identification of angels with pre-existent or disembodied human souls, a confusion later condemned in its Origenist form at Constantinople (543/553) and rejected by Augustine and St.~John of Damascus: angelic nature and human nature are distinct created species. Second, following Psalm~95:5 LXX, the Church rejected Philo's neutral equation of biblical angels with philosophical \gk{daimones} (Augustine, \emph{City of God} IX.19). Third, against any notion of a sub-divine Logos or of angels assisting God to fabricate human bodies or lower matter, St.~Irenaeus insisted that the Father had no need of angels to call creation into being or to form man, since He creates directly through His Son and Spirit, whom all angels serve (\emph{Against Heresies} II.30.6, IV.7.4, IV.20.1), a principle reaffirmed by St.~Basil (\emph{Hexaemeron}) and St.~Augustine (\emph{On the Trinity}~III).

\subsection{Greek Poets and Philosophers Under Patristic Scrutiny}
\label{subsec:greek-philosophers}

Classical and late-antique Greek thought possessed a dense reflection on intermediate spirits long before the Gospel reached Athens. Hesiod's \emph{Works and Days} recounted how men of the golden race became guardian \gk{daimones} roaming the earth. Plato gave that middle realm philosophical form in the \emph{Symposium} (202d--203a), where Diotima classifies \gk{Er\={o}s} as a great \gk{daim\={o}n} occupying the space between the divine and the mortal (\gk{metaxy}), mediating entreaties, sacrifices, divination, priestcraft, incantations, and sorcery; in the \emph{Timaeus}, Plato spoke of lesser immortal ``gods'' created by the supreme Demiurge. Middle Platonists such as Plutarch (\emph{De defectu oraculorum}, \emph{De genio Socratis}) and Apuleius (\emph{De deo Socratis}) systematized these \gk{daimones} as aerial rational animals bridging heaven and earth. Later Neoplatonists---Plotinus (\emph{Enneads}), Porphyry (\emph{Letter to Anebo}, \emph{De regressu animae}), and Iamblichus (\emph{De mysteriis})---integrated intermediary spirits into emanated triadic hierarchies and, especially in Iamblichus, into ``theurgy'' (\gk{theourgia}), ritual invocation designed to purify the soul through spiritual powers.

St.~Augustine subjected this entire tradition to a searching, sustained examination in Books~VIII through~X of \emph{The City of God}, quoting its foremost spokesmen directly and measuring every claim against Scripture and the mystery of Christ. First, he dismantles the Platonist hierarchy of aerial mediators (\emph{City of God} VIII.14--22; IX.15). Summarizing their system, he writes:
\begin{quotation}
There is, say they, a threefold division of all animals endowed with a rational soul, namely, into gods, men, and demons. The gods occupy the loftiest region, men the lowest, the demons the middle region. For the abode of the gods is heaven, that of men the earth, that of the demons the air... For they have immortality of body in common with the gods, but passions of the mind in common with men. (St.~Augustine, \emph{City of God} VIII.14)
\end{quotation}
Two chapters later, Augustine preserves Apuleius's exact fivefold definition from \emph{De deo Socratis}: ``Demons are of an animal nature, passive in soul, rational in mind, aerial in body, eternal in time'' (\emph{City of God} VIII.16). From that very definition Augustine springs his refutation. A spirit enslaved to anger, lust, envy, and theatrical obscenity while locked forever in an immortal body is not superior to a virtuous mortal man, but far more wretched. Even the Neoplatonist Plotinus had unwittingly exposed their misery:
\begin{quotation}
In speaking of human souls, he says, ``The Father in compassion made their bonds mortal;'' that is to say, he considered it due to the Father's mercy that men, having a mortal body, should not be for ever confined in the misery of this life. But of this mercy the demons have been judged unworthy, and they have received, in conjunction with a soul subject to passions, a body not mortal like man's, but eternal. (St.~Augustine, \emph{City of God} IX.10)
\end{quotation}
Because passion-ridden spirits cannot reconcile sinners to a holy God, the true Mediator must share immortality with God and sinless mortality with men: Jesus Christ alone stands as Mediator between God and humanity (1~Tim~2:5; \emph{City of God} IX.15).

Second, Augustine settles the Christian vocabulary once for all against the Platonic category of benevolent demons:
\begin{quotation}
The Platonists do not deny their existence, but prefer to call them good demons. But we, following Scripture, according to which we are Christians, have learned that some of the angels are good, some bad, but never have we read in Scripture of good demons; but wherever this or any cognate term occurs, it is applied only to wicked spirits. (St.~Augustine, \emph{City of God} IX.19)
\end{quotation}
Only after closing the door on ``good demons'' does he extend the serene generosity of Book~IX, chapter~23 toward Plato's \emph{Timaeus}: if philosophers acknowledge immortal spirits created by the supreme God and blessed solely by cleaving to Him, they confess the reality of the holy angels, whatever name they employ.

Third, in Book~X, Augustine turns to Porphyry and Iamblichean theurgy. He strips away the pretension that ``theurgy'' is a nobler art than vulgar magic, noting that ``both classes are the slaves of the deceitful rites of the demons whom they invoke under the names of angels. For even Porphyry promises some kind of purgation of the soul by the help of theurgy, though he does so with some hesitation and shame'' (\emph{City of God} X.9). While commending Porphyry's \emph{Letter to Anebo} for exposing sacrifice-hungry spirits (\emph{City of God} X.11), Augustine faults his timidity in still calling certain demons benignant. Yet Augustine gladly receives one sterling maxim from Porphyry himself:
\begin{quotation}
And hence, even this Platonist himself judiciously observes that we should rather imitate than invoke them. We ought not, then, to fear that we may offend these immortal and happy subjects of the one God by not sacrificing to them; for this they know to be due only to the one true God, in allegiance to whom they themselves find their blessedness. (St.~Augustine, \emph{City of God} X.26)
\end{quotation}
Holy angels desire our fellowship in adoring God, never sacrifices offered to themselves.

A century after Augustine, Pseudo-Dionysius completed the patristic reception of Neoplatonism from the East. In \emph{The Celestial Hierarchy} (chs.~1--4) and \emph{On the Divine Names} (ch.~4), he received the Neoplatonic language of triadic order, illumination, and graded participation in the Good, yet inverted its metaphysical foundation at every decisive point. Angels do not emanate necessarily from the One, but are created freely out of nothing by the Triune God; every hierarchy has Jesus Christ as its sole origin, leader, and deifying light (\CH{1.2}, \CH{4.4}); ritual invocation of subordinate gods vanishes into the sacramental liturgy of the Church; and against any dualism of evil spirits, Dionysius establishes from the metaphysics of goodness that ``nor are the demons evil by nature'' (\emph{On the Divine Names} IV.23), since every created nature is good and demonic malice is solely a voluntary falling away from the Good.
